\section{A nonabelian kernel for every parameter}
% BEGIN F1.Construction
\begin{theorem}[The cyclotomic polyhedra]\label{F1.Construction}
For every odd prime $p$ and every $m\geq3$, there is a string
C-group of order $4p^m$ and type $\{p,2p\}$ with nonabelian
Sylow $p$-subgroup. More precisely, let
\[
 O=\mathbb Z[\zeta_p],\quad \pi=\zeta_p-1,\quad
 A=O/\pi^{m-1}O,\quad P=A\rtimes\gen{\beta},
 \qquad |\beta|=p,\quad v^\beta=\zeta_pv.
\]
Define automorphisms by
\[
 \sigma(v)=\overline v,\quad\sigma(\beta)=\beta^{-1},
 \qquad\tau(v)=-v,\quad\tau(\beta)=\beta.
\]
Then $G=P\rtimes\gen{\sigma,\tau}$ has the required marking
\[
 r_0=s\tau\sigma,\qquad r_1=\beta\tau\sigma,\qquad r_2=\sigma,
\]
where $s$ denotes the additive class of the ring element $1$.
It is not the identity of the additive group $A$.
\end{theorem}
\begin{proof}
Set $n=m-1\geq2$ and write $\zeta=\zeta_p$. To fix the action
convention, use pairs $(v,k)\in A\times\mathbb Z/p\mathbb Z$ with
\begin{equation}\label{eq:pairlaw}
 (v,k)(w,l)=(v+\zeta^{-k}w,k+l).
\end{equation}
Then $\beta=(0,1)$ has $v^\beta=\zeta v$.
The cyclotomic coordinates of~\cite{UF} give $|P|=p^{n+1}=p^m$.
The action by $\zeta$ is nontrivial for $n\geq2$, so $P$ is
nonabelian. Every element outside $A$ has order exactly $p$:
for $k\ne0$ its $p$th power has first coordinate
$\sum_{i=0}^{p-1}\zeta^{-ik}v=0$, and its image in $P/A$ has
order $p$.

On pairs the proposed maps are
$\sigma(v,k)=(\overline v,-k)$ and $\tau(v,k)=(-v,k)$.
They preserve~\eqref{eq:pairlaw}, since conjugation sends
$\zeta^{-k}$ to $\zeta^k$. They are commuting involutory
automorphisms. Moreover, $\sigma$ inverts the nontrivial element
$\beta$, whereas $\tau$ fixes it; and $\tau$ is nontrivial on
the nonzero odd-order group $A$. Thus they are independent and
generate a Klein four group. In particular $|G|=4p^m$.

Put $h=\tau\sigma$. In multiplicative group notation
$s^h=s^{-1}$ and $\beta^h=\beta^{-1}$. The elementary reflection
identities of~\cite{UF} make $r_0=sh$, $r_1=\beta h$ and
$r_2=\sigma$ involutions. Directly,
\[
 r_0r_2=s\tau=r_2r_0,\qquad
 r_0r_1=s\beta^{-1},\qquad r_1r_2=\beta\tau.
\]
The middle product is outside $A$ and has order $p$.
The last has order $2p$, since its two factors commute and have
coprime orders. The marked generators are nonidentity and distinct;
the first two project to $\tau\sigma$ in the complement, the third
to $\sigma$, and the first two are distinct because $s\ne\beta$.

They generate the whole group. Indeed,
$(r_1r_2)^p=\tau$ yields $\tau$, then $\beta$; the generator
$r_2$ gives $\sigma$, and $r_0$ gives $s$. The
$\beta$-conjugates of $s$ are $1,\zeta,\ldots,\zeta^{p-1}$,
whose classes additively generate $A$.

Finally $H=\gen{r_0,r_1}$ and $J=\gen{r_1,r_2}$ have
orders $2p$ and $4p$. The subgroup $C=\gen{r_1}$ has prime
index $p$ in $H$. If $H\leq J$, all three generators would lie
in $J$, giving $G=J$ and contradicting $4p^m>4p$.
The prime-index lemma of~\cite{UF} therefore gives $H\cap J=C$.
Its rank-three test proves the full intersection property.
\end{proof}
% END F1.Construction
